\chapter{Apr.~21 --- Surfaces}

\section{Intersection Numbers}

\begin{remark}
  Recall that if $X$ is a projective
  variety and $\mathcal{F} \in \Coh_X$,
  then
  \[
    \chi(X, \mathcal{F} \otimes \OO_X(m))
  \]
  (the \emph{Hilbert polynomial})
  is a polynomial of degree $\dim \Supp(\mathcal{F})$
  in $m$.
  When $\dim X = 1$, we have an explicit
  description (by Riemann-Roch) as follows:
  \begin{quote}
    \vspace{-2em}
    \begin{theorem}
      If $C$ is a curve and $D$ is a
      Cartier divisor on $C$, then
      \[
        \chi(C, \OO_C(mD))
        = m \deg D + 1 - g(C).
      \]
    \end{theorem}
  \end{quote}
  How does this generalize to higher
  dimensions?

  Let $X \subseteq \PP^n$ be a smooth
  projective surface and $D$ a divisor on
  $X$. We want to extract something like
  ``$\deg D$.'' The idea is the following:
  Take a general hyperplane 
  $H \subseteq \PP^n$, and consider
  $\#(D \cap H)$. Note
  \[
    D \cap H
    = D \cap (H \cap X),
  \]
  where $D$ and $H \cap X$ are both curves.
\end{remark}

\begin{theorem}
  There exists a symmetric bilinear pairing
  \begin{align*}
    \Div X \times \Div X
    &\longrightarrow \Z \\
    (C, D) &\longmapsto C \cdot D
  \end{align*}
  such that
  \begin{enumerate}
    \item (linear equivalence) if
      $C, D, C', D' \in \Div X$ with
      $C \sim C'$ and $D \sim D'$, then
      $C \cdot D = C' \cdot D'$;
    \item if $C, D$ are distinct irreducible
      curves, then
      $C \cdot D = \sum_{p \in C \cap D} (C \cdot D)_p$,
      where
      \[
        (C \cdot D)_p
        = \text{local intersection number at $p$}
        = \dim_k (\OO_{X, p} / (f, g))
      \]
      with $I_{C, p} = (f)$ and
      $I_{D, p} = (g)$.
  \end{enumerate}
\end{theorem}

\begin{remark}
  If $C$ and $D$ are smooth
  curves meeting transversely at
  $p$ (i.e. the tangent directions of
  $C, D$ at $p$ are
  distinct), then
  $(C \cdot D)_p = 1$.

  To see this, note that $C, D$ meet
  transversely at $p$ if and only if
  $k[f] \ne k[g]$ in
  $\m_p / \m_p^2$, which happens if and only
  if $[f], [g]$ generate
  $\m_p / \m_p^2$. By Nakayama's lemma,
  this happens if and only if
  $f, g$ generate $\m_p$, i.e.
  $(f, g) = \m_p$. This implies that
  $\dim_k (\OO_{X, p} / (f, g)) = \dim_k k = 1$.
\end{remark}

\begin{remark}
  The pairing $\Div X \times \Div X \to \Z$
  induces a symmetric bilinear
  form
  \begin{align*}
    \Cl(X) \times \Cl(X) &\longrightarrow \Z \\
    ([C], [D]) &\longmapsto C \cdot D.
  \end{align*}
\end{remark}

\begin{example}
  Consider $\PP^2$, where
  $\Cl(\PP^2) = \Z[H]$ for some line
  $H \subseteq \PP^2$.
  Then
  \[
    H \cdot H
    = H \cdot H'
    = 1,
  \]
  where $H'$ is a line different
  from $H$ in $\PP^2$.
  So if $C, D \subseteq \PP^2$ are
  plane curves of degrees $c, d$, then
  \[
    C \cdot D
    = (cH) \cdot (dH)
    = cd H^2 = cd.
  \]
  This recovers the version of
  Bezout's theorem which
  counts multiplicity.
\end{example}

\begin{example}
  Consider $\PP^1_{x_0 : x_1} \times \PP^1_{y_0 : y_1}$,
  where $\Cl(\PP^1 \times \PP^1) = \Z[H_1] \oplus \Z[H_2]$
  for
  \[
    H_1 = \{[0 : 1]\} \times \PP^1
    \quad \text{and} \quad
    H_2 = \PP^1 \times \{[0 : 1]\}.
  \]
  Then $H_1 \cdot H_1 = H_1 \cdot ([1 : 0] \times \PP^1) = 0$
  (similarly, $H_2 \cdot H_2 = 0$) and
  $H_1 \cdot H_2 = 1$.
\end{example}

\begin{example}
  Let
  $\pi : B_p \PP^2 \to \PP^2$ be the
  blowup and
  $E \subseteq B_p \PP^2$ the
  exceptional divisor. Then
  \[
    \Cl(B_p \PP^2)
    = \Z[E] \oplus \Z[\widetilde{H}],
  \]
  where $p \notin H \subseteq \PP^2$ is a
  hyperplane. Here $\widetilde{H}$ denotes
  the strict transform. Then we can compute:
  \begin{enumerate}
    \item $\widetilde{H} \cdot \widetilde{H}$. Choose
      a line $p \notin H' \subseteq \PP^2$
      different from $H$. Now
      $H \sim H'$, so
      \[
        \widetilde{H}
        = \pi^* H
        \sim \pi^* H'
        = \widetilde{H}',
      \]
      and thus $\widetilde{H} \cdot \widetilde{H} = \widetilde{H} \cdot \widetilde{H}' = 1$.
    \item $E \cdot \widetilde{H} = 0$.
    \item $E \cdot E$. Choose distinct
      lines $p \in L_i \subseteq \PP^2$ for
      $i = 1, 2$. Now
      $L_1 \sim L_2 \sim H$, and
      \[
        \pi^* L_i
        = E + \widetilde{L}_i.
      \]
      So
      $E \sim \pi^* L_i - \widetilde{L}_i \sim \widetilde{H} - \widetilde{L}_i$,
      and thus
      \begin{align*}
        E \cdot E
        = (\widetilde{H} - \widetilde{L}_1)
        \cdot (\widetilde{H} - \widetilde{L}_2)
        &= (\widetilde{H} \cdot \widetilde{H})
        - (\widetilde{L}_1 \cdot \widetilde{H})
        - (\widetilde{L}_2 \cdot \widetilde{H})
        + (\widetilde{L}_1 \cdot \widetilde{L}_2) \\
        &= 1
        - (\widetilde{L}_1 \cdot \widetilde{H})
        - (\widetilde{L}_2 \cdot \widetilde{H})
        + (\widetilde{L}_1 \cdot \widetilde{L}_2).
      \end{align*}
      Note that $\widetilde{L}_1$ and
      $\widetilde{L}_2$ do not intersect in
      the blowup, while
      $\widetilde{L}_i$ and $\widetilde{H}$
      intersect at a point. Thus
      \[
        E \cdot E
        = 1 - 1 - 1 + 0 = -1.
      \]
  \end{enumerate}
\end{example}

\begin{lemma}
  If $C$ and $D$ are distinct curves
  on $X$ with $C$ smooth and irreducible, then
  \[
    C \cdot D
    = \deg(\OO_X(D)|_C).
  \]
\end{lemma}

\begin{proof}
  We have a short exact sequence
  \[
    \begin{tikzcd}
      0 \ar[r] & \mathcal{I}_D
      \ar[r] & \OO_X
      \ar[r] & \OO_D
      \ar[r] & 0
    \end{tikzcd}
  \]
  and applying $- \otimes \OO_C$ gives a
  short exact sequence
  (recall that $\OO_D = \OO_X / \mathcal{I}_D$)
  \[
    \begin{tikzcd}
      0 \ar[r] &
      \OO_X(-D) \otimes \OO_C
      \ar[r] & \OO_C
      \ar[r] & \OO_X / (\mathcal{I}_C + \mathcal{I}_D)
      \ar[r] & 0
    \end{tikzcd}
  \]
  Thus the Euler characteristics satisfy
  \[
    \chi(C, \OO_C)
    - \chi(C, \OO_X(-D)|_C)
    = \chi(X, \OO_X / (\mathcal{I}_C + \mathcal{I}_D))
    = h^0(X, \OO_X / (\mathcal{I}_C + \mathcal{I}_D))
    = C \cdot D,
  \]
  since $\OO_X / (\mathcal{I}_C + \mathcal{I}_D)$ is supported at finitely many points,
  so $h^1(X, \OO_X / (\mathcal{I}_C + \mathcal{I}_D)) = 0$.
  Also,
  \[
    \chi(C, \OO_C)
    - \chi(C, \OO_X(-D)|_C)
    = 0 - \deg(\OO_X(-D)|_C)
    = \deg(\OO_X(D)|_C)
  \]
  by the Riemann-Roch theorem, which
  gives the desired formula.
\end{proof}

\begin{example}
  If $C$ is a smooth curve on $X$, then
  $\omega_C
    = \omega_X(C)|_C$, where
  $\omega_X(C) = \OO_X(K_X + C)$.
  This implies that
  \[
    2 g(C) - 2
    = \deg \omega_C
    = \deg(\omega_X(C)|_C)
    = (K_X + C) \cdot C.
  \]
  As a special case, if $C$ is a smooth
  degree $d$ plane curve in $\PP^2$, then
  we see that
  \[
    2g(C) - 2
    = (K_{\PP^2} + C) \cdot C
    = (-3H + d H) \cdot dH
    = (d - 3)d,
  \]
  where $H \subseteq \PP^2$ is a line.
  This recovers the formula
  $g(C) = (d - 1)(d - 2) / 2$.
\end{example}

\section{Riemann-Roch for Surfaces}

\begin{theorem}[Riemann-Roch for surfaces]\label{thm:RR-surfaces}
  For a divisor $D$ on a smooth
  projective surface $X$,
  \[
    \chi(X, \OO_X(D))
    = \frac{1}{2} D \cdot (D - K_X)
    + 1 + p_a(X),
  \]
  where $p_a(X) = \chi(X, \OO_X) - 1$
  is the \emph{arithmetic genus} of $X$.
\end{theorem}

\begin{lemma}\label{lem:divisor-decomposition}
  Any divisor $D$ on a smooth
  projective surface $X$ can be written
  as $D \sim C' - C$ with
  $C$ and $C'$ smooth irreducible curves.
\end{lemma}

\begin{proof}
  Choose a very ample divisor $H$ on $X$.
  For $k \gg 0$, the divisor
  $D + kH$ is very ample.
  By Bertini's theorem (see
  Mustata 6.4.1), there exist
  smooth irreducible curves
  $C \sim kH$ and $C' \sim D + kH$.
  Then we can write
  $D = (D + kH) - kH \sim C' - C$.
\end{proof}

\begin{proof}[Proof of Theorem \ref{thm:RR-surfaces}]
  The goal is to reduce to the case
  of curves.
  Since the left and right hand sides
  only depend on the linear equivalence class
  of $D$, we may assume that $D = C - E$ 
  with $C, E$ smooth irreducible curves
  (by Lemma \ref{lem:divisor-decomposition}).
  We have a short exact sequence
  \[
    \begin{tikzcd}
      0 \ar[r] & \OO_X(C - E)
      \ar[r] & \OO_X(C) \ar[r] & \OO_X(C) \otimes \OO_E
      \ar[r] & 0
    \end{tikzcd}
  \]
  We also have a second short exact sequence
  \[
    \begin{tikzcd}
      0 \ar[r]
      & \OO_X
      \ar[r] & \OO_X(C)
      \ar[r] & \OO_X(C) \otimes \OO_C
      \ar[r] & 0
    \end{tikzcd}
  \]
  \pagebreak
  Thus taking Euler characteristics, we see
  that
  \begin{align*}
    \chi(X, \OO_X(C - E))
    &= \chi(X, \OO_X(C))
    - \chi(X, \OO_X(C) \otimes \OO_E) \\
    &= \chi(X, \OO_X)
    + \chi(X, \OO_X(C) \otimes \OO_C)
    - \chi(X, \OO_X(C) \otimes \OO_E) \\
    &= \chi(X, \OO_X)
    + \chi(C, \OO_X(C)|_C)
    - \chi(E, \OO_X(C)|_E) \\
    &= (1 + p_a(X))
    + (C^2 + 1 - g(C))
    - (C \cdot E + 1 - g(E)).
  \end{align*}
  Now
  $g(C) = \frac{1}{2} C \cdot (C + K_X) + 1$
  and $g(E) = \frac{1}{2} E \cdot (E + K_X) + 1$,
  so we get
  \[
    \chi(X, \OO_X(C - E))
    = \frac{1}{2}(C - E) \cdot (C - E - K_X)
    + 1 + p_a(X).
  \]
  Since $D = C - E$, this gives the
  desired formula.
\end{proof}

\begin{remark}
  There is a criterion for ampleness
  in terms of intersection numbers
  (the \emph{Nakai-Moishezon criterion}):
  A divisor $D$ on $X$ is ample
  if and only if $D^2 > 0$ and
  $D \cdot C > 0$ for all irreducible
  curves $C \subseteq X$. This is
  generally true in any dimension.
\end{remark}
